% Table: Placebo regressions (Chapter 4 appendix)
\begin{table}[htbp]
\centering
\caption{Placebo test: IV estimates for students who do not take A-level economics}
\label{tab:placebo}
\begin{tabular}{l cc}
\toprule
                                      & (1)            & (2)            \\
                                      & Full sample    & Non-A-level    \\
                                      &                & econ students  \\
\midrule
\multicolumn{3}{c}{\textbf{Panel A: Log earnings at age 28}} \\
\midrule
A-level econ available                & 0.008\sym{*}    & -0.002             \\
                                      & (0.004)           & (0.004)           \\
\addlinespace
School FE                             & \checkmark     & \checkmark     \\
Cohort FE                             & \checkmark     & \checkmark     \\
Individual controls                   & \checkmark     & \checkmark     \\
\addlinespace
Observations                          & 368,260             & 357,168             \\
\bottomrule
\end{tabular}
\begin{minipage}{\textwidth}
\smallskip
\footnotesize
\textit{Notes:} Column~(1) reproduces the reduced-form estimates from the full sample for reference. Column~(2) restricts the sample to students who do not take A-level economics. If availability affects earnings for non-takers, this would indicate a direct effect of availability on outcomes, violating the exclusion restriction. Because availability itself moves students out of the non-taker group, the test conditions on a post-treatment choice and is supportive rather than dispositive. Standard errors clustered at the school level in parentheses. Sample: students in switcher schools, GCSE cohorts 2002--2008. \sym{*} $p<0.10$, \sym{**} $p<0.05$, \sym{***} $p<0.01$.
\end{minipage}
\end{table}
